\documentclass[a4paper,11pt]{article}
\usepackage[utf8]{inputenc}
\usepackage[T1]{fontenc}
\usepackage[french]{babel}
\usepackage{booktabs}
\usepackage{amsmath}
\usepackage{geometry}
\geometry{margin=2.5cm}

\title{Prosit 2 -- Machine Learning\\{\large De la donn\'ee brute au Perceptron Multi-Couches}}
\author{Prosit 2 -- Machine Learning}
\date{Mars 2026}

\begin{document}
\maketitle
\tableofcontents
\newpage

% ============================================================
\part{Pr\'eparation des donn\'ees}
% ============================================================

\section{Objectif}

Transformer le dataset brut \texttt{housing.csv} (20\,640 lignes, 10 colonnes) en un fichier pr\^et pour l'entra\^inement d'un perceptron : \texttt{amputed\_normalized\_housing\_data.csv}.

Pipeline : imputation KNN $\to$ encodage ordinal $\to$ nettoyage outliers $\to$ normalisation Z-Score.

\section{\'Etape 1 -- Imputation KNN}

La colonne \texttt{total\_bedrooms} contient 207 valeurs manquantes. On les comble par KNN ($k=5$, pond\'eration par distance inverse) :

\begin{equation}
\hat{x}_i = \frac{\sum_{j \in \mathcal{N}_k(i)} w_j \cdot x_j}{\sum_{j \in \mathcal{N}_k(i)} w_j}, \quad w_j = \frac{1}{d(i,j)}
\end{equation}

Les donn\'ees sont standardis\'ees temporairement avant KNN pour que la distance euclidienne soit coh\'erente entre colonnes d'\'echelles diff\'erentes, puis d\'e-standardis\'ees apr\`es imputation.

\section{\'Etape 2 -- Encodage ordinal de \texttt{ocean\_proximity}}

\begin{table}[h]
\centering
\begin{tabular}{lc}
\toprule
\textbf{Cat\'egorie} & \textbf{Valeur} \\
\midrule
\texttt{<1H OCEAN}  & 0 \\
\texttt{INLAND}     & 1 \\
\texttt{ISLAND}     & 2 \\
\texttt{NEAR BAY}   & 3 \\
\texttt{NEAR OCEAN} & 4 \\
\bottomrule
\end{tabular}
\caption{Encodage ordinal de \texttt{ocean\_proximity}}
\end{table}

Cet encodage est r\'ealis\'e \textbf{avant} l'imputation KNN (KNNImputer requiert des donn\'ees num\'eriques).

\section{\'Etape 3 -- Nettoyage des outliers}

Deux types de nettoyage sont appliqu\'es apr\`es imputation :

\subsection{Suppression des prix plafonn\'es}

965 lignes avec \texttt{median\_house\_value} $= 500\,001$\,\$ sont supprim\'ees. Ces valeurs sont \textbf{censur\'ees} (le prix r\'eel est sup\'erieur mais non enregistr\'e), ce qui biaiserait l'apprentissage.

\subsection{Suppression des outliers IQR}

Pour chaque feature num\'erique (sauf \texttt{ocean\_proximity} et \texttt{median\_house\_value}), on supprime les lignes dont la valeur est en dehors de l'intervalle :

\begin{equation}
[Q_1 - 1.5 \times IQR,\; Q_3 + 1.5 \times IQR]
\end{equation}

o\`u $IQR = Q_3 - Q_1$ est l'\'ecart interquartile.

\begin{table}[h]
\centering
\begin{tabular}{lrc}
\toprule
\textbf{Feature} & \textbf{Outliers} & \textbf{Bornes} \\
\midrule
\texttt{total\_rooms}    & 1\,251 & $[-1085, 5643]$ \\
\texttt{total\_bedrooms} & 1\,227 & $[-230, 1175]$ \\
\texttt{population}      & 1\,126 & $[-629, 3171]$ \\
\texttt{households}      & 1\,165 & $[-204, 1092]$ \\
\texttt{median\_income}  & 364    & $[-0.6, 7.7]$ \\
\bottomrule
\end{tabular}
\caption{Outliers d\'etect\'es par feature (m\'ethode IQR)}
\end{table}

\textbf{Bilan} : 3\,039 lignes supprim\'ees (14.7\%), \textbf{17\,601 lignes restantes}.

\section{\'Etape 4 -- Normalisation Z-Score}

Toutes les features sauf \texttt{median\_house\_value} (variable cible) sont normalis\'ees :

\begin{equation}
z = \frac{x - \mu}{\sigma}
\end{equation}

\textbf{Justification} : le perceptron calcule $\hat{y} = f(\mathbf{w}^T \mathbf{x} + b)$. Sans normalisation, les features \`a grande \'echelle (\texttt{total\_rooms} $\sim 2000$) domineraient le gradient face aux petites (\texttt{median\_income} $\sim 4$).

\texttt{median\_house\_value} n'est pas transform\'ee car c'est la cible de la r\'egression.

\section{Sch\'ema du dataset final}

\begin{table}[h]
\centering
\begin{tabular}{lc}
\toprule
\textbf{Colonne} & \textbf{Transformation} \\
\midrule
\texttt{longitude}            & Z-Score \\
\texttt{latitude}             & Z-Score \\
\texttt{housing\_median\_age} & Z-Score \\
\texttt{total\_rooms}         & Z-Score \\
\texttt{total\_bedrooms}      & KNN + Z-Score \\
\texttt{population}           & Z-Score \\
\texttt{households}           & Z-Score \\
\texttt{median\_income}       & Z-Score \\
\texttt{median\_house\_value} & \textbf{Non (cible)} \\
\texttt{ocean\_proximity}     & Ordinal + Z-Score \\
\bottomrule
\end{tabular}
\caption{Sch\'ema de \texttt{amputed\_normalized\_housing\_data.csv} -- 17\,601 lignes, 10 colonnes}
\end{table}

% ============================================================
\part{Mod\'elisation -- Perceptron Multi-Couches}
% ============================================================

\section{Analyse de colin\'earit\'e}

Avant l'entra\^inement, on supprime les features redondantes (corr\'elation $|r| > 0.8$). Pour chaque paire colin\'eaire, on conserve la feature la mieux corr\'el\'ee avec la cible.

\begin{table}[h]
\centering
\begin{tabular}{llcl}
\toprule
\textbf{Feature A} & \textbf{Feature B} & $|r|$ & \textbf{Supprim\'ee} \\
\midrule
longitude      & latitude        & 0.925 & longitude \\
total\_rooms   & total\_bedrooms & 0.930 & total\_bedrooms \\
total\_rooms   & population      & 0.857 & population \\
total\_rooms   & households      & 0.918 & households \\
total\_bedrooms & population     & 0.878 & population \\
total\_bedrooms & households     & 0.980 & total\_bedrooms \\
population     & households      & 0.907 & population \\
\bottomrule
\end{tabular}
\caption{Paires colin\'eaires et features supprim\'ees}
\end{table}

\textbf{Features conserv\'ees (5)} : \texttt{latitude}, \texttt{housing\_median\_age}, \texttt{total\_rooms}, \texttt{median\_income}, \texttt{ocean\_proximity}.

\textbf{Features supprim\'ees (4)} : \texttt{longitude}, \texttt{total\_bedrooms}, \texttt{population}, \texttt{households}.

\section{S\'eparation Train / Test}

Le dataset (17\,601 lignes apr\`es nettoyage) est divis\'e al\'eatoirement (\texttt{random\_state=42}) :

\begin{itemize}
    \item \textbf{Train} : 12\,320 \'echantillons (70\%)
    \item \textbf{Test} : 5\,281 \'echantillons (30\%)
\end{itemize}

\section{Architecture du MLP}

Perceptron multi-couches \`a 2 couches cach\'ees :

\begin{center}
\textbf{5} (entr\'ee) $\to$ \textbf{64} (ReLU) $\to$ \textbf{32} (ReLU) $\to$ \textbf{1} (sortie, identit\'e)
\end{center}

Propagation avant :
\begin{align}
h_1 &= \text{ReLU}(W_1 x + b_1) \\
h_2 &= \text{ReLU}(W_2 h_1 + b_2) \\
\hat{y} &= W_3 h_2 + b_3
\end{align}

O\`u $\text{ReLU}(z) = \max(0, z)$.

\textbf{Hyperparam\`etres} :
\begin{itemize}
    \item Optimiseur : Adam (learning rate = 0.001)
    \item Loss : MSE (Mean Squared Error)
    \item Early stopping sur 10\% de validation
    \item Max 500 epochs --- converg\'e en \textbf{271 epochs}
    \item Nombre total de poids : \textbf{2\,497}
\end{itemize}

% ============================================================
\part{\'Evaluation du mod\`ele}
% ============================================================

\section{M\'etriques utilis\'ees}

C'est un probl\`eme de \textbf{r\'egression}, pas de classification : il n'y a pas de << pr\'ecision >> (\textit{accuracy}) au sens classique. On utilise 4 m\'etriques compl\'ementaires.

\subsection{MAE -- Mean Absolute Error}

\begin{equation}
\text{MAE} = \frac{1}{n}\sum_{i=1}^{n} |y_i - \hat{y}_i|
\end{equation}

\textbf{\'Echelle} : dollars (\$). Erreur moyenne en valeur absolue --- robuste aux outliers.

\subsection{RMSE -- Root Mean Squared Error}

\begin{equation}
\text{RMSE} = \sqrt{\frac{1}{n}\sum_{i=1}^{n} (y_i - \hat{y}_i)^2}
\end{equation}

\textbf{\'Echelle} : dollars (\$). P\'enalise les grosses erreurs plus fortement que la MAE (effet quadratique).

\subsection{MAPE -- Mean Absolute Percentage Error}

\begin{equation}
\text{MAPE} = \frac{100}{n}\sum_{i=1}^{n} \left|\frac{y_i - \hat{y}_i}{y_i}\right|
\end{equation}

\textbf{\'Echelle} : pourcentage (\%). Erreur relative --- ind\'ependante de l'\'echelle des prix. On peut l'interpr\'eter comme : << en moyenne, le mod\`ele se trompe de X\% du prix r\'eel >>.

\subsection{R\textsuperscript{2} -- Coefficient de d\'etermination}

\begin{equation}
R^2 = 1 - \frac{\sum_{i=1}^{n}(y_i - \hat{y}_i)^2}{\sum_{i=1}^{n}(y_i - \bar{y})^2}
\end{equation}

\textbf{\'Echelle} : sans unit\'e, entre 0 et 1. Part de la variance du prix expliqu\'ee par le mod\`ele. C'est la m\'etrique la plus proche d'une << pr\'ecision >> pour la r\'egression.

\section{R\'esultats}

\subsection{Avant nettoyage (20\,640 lignes)}

\begin{table}[h]
\centering
\begin{tabular}{lcccc}
\toprule
& \textbf{MAE (\$)} & \textbf{RMSE (\$)} & \textbf{MAPE (\%)} & \textbf{R\textsuperscript{2}} \\
\midrule
Train (70\%) & 52\,455 & 73\,146 & 29.85 & 0.60 \\
Test (30\%)  & 52\,183 & 72\,690 & 29.97 & 0.60 \\
\bottomrule
\end{tabular}
\caption{Performance avant nettoyage des outliers}
\end{table}

\subsection{Apr\`es nettoyage (17\,601 lignes)}

\begin{table}[h]
\centering
\begin{tabular}{lcccc}
\toprule
& \textbf{MAE (\$)} & \textbf{RMSE (\$)} & \textbf{MAPE (\%)} & \textbf{R\textsuperscript{2}} \\
\midrule
Train (70\%) & 45\,840 & 62\,841 & 28.24 & 0.56 \\
Test (30\%)  & 46\,288 & 63\,958 & 27.98 & 0.55 \\
\bottomrule
\end{tabular}
\caption{Performance apr\`es nettoyage des outliers}
\end{table}

\subsection{Comparaison}

\begin{table}[h]
\centering
\begin{tabular}{lccc}
\toprule
\textbf{M\'etrique (test)} & \textbf{Avant} & \textbf{Apr\`es} & \textbf{Am\'elioration} \\
\midrule
MAE (\$)  & 52\,183 & 46\,288 & $-$11.3\% \\
RMSE (\$) & 72\,690 & 63\,958 & $-$12.0\% \\
MAPE (\%) & 29.97   & 27.98   & $-$2.0 pts \\
R\textsuperscript{2} & 0.60 & 0.55 & $-$0.05 \\
\bottomrule
\end{tabular}
\caption{Impact du nettoyage des outliers sur la performance}
\end{table}

\section{Comment lire ces r\'esultats ?}

\begin{itemize}
    \item \textbf{MAE et RMSE en baisse} : le mod\`ele fait des erreurs plus petites en valeur absolue ($-$12\% de RMSE). C'est l'am\'elioration la plus concr\`ete.
    \item \textbf{MAPE en baisse (28\%)} : en moyenne, la pr\'ediction s'\'ecarte de 28\% du prix r\'eel, soit une << pr\'ecision relative >> d'environ \textbf{72\%}.
    \item \textbf{R\textsuperscript{2} l\'eg\`erement plus bas (0.55 vs 0.60)} : normal --- en supprimant les outliers, on r\'eduit la variance totale des prix. Le d\'enominateur de $R^2$ diminue, donc le ratio peut baisser m\^eme si les erreurs absolues sont plus petites. La MAE et le RMSE sont de meilleurs indicateurs ici.
    \item \textbf{Train $\approx$ Test} : pas de sur-apprentissage.
\end{itemize}

\section{Limites identifi\'ees}

\begin{itemize}
    \item 4 features supprim\'ees par colin\'earit\'e --- un mod\`ele non-lin\'eaire pourrait en tirer plus d'information.
    \item Seules 5 features utilis\'ees --- des donn\'ees suppl\'ementaires (surface, nombre de pi\`eces, quartier) am\'elioreraient la pr\'ediction.
    \item Le seuil IQR ($1.5 \times$) est un choix conservateur ; un seuil plus souple ($2 \times$ ou $3 \times$) garderait plus de donn\'ees.
\end{itemize}

\end{document}
